\documentclass[11pt,a4paper]{article}

\usepackage[utf8]{inputenc}
\usepackage{amsmath,amssymb,amsfonts}
\usepackage{geometry}
\geometry{margin=1in}
\usepackage{booktabs}
\usepackage{hyperref}
\usepackage{cite}
\usepackage{microtype}
\usepackage{setspace}
\onehalfspacing

\title{\textbf{Title of the Working Paper: \\ \large Subtitle or Focus Area}}

\author{
    \textbf{Author Name or Desk}\thanks{Correspondence: \texttt{research@seekkey.ltd}. Sovereign corporate entity: Seek Key Ltd (United Kingdom). This working paper is part of the Seek Key Epistemological Archive. Hypotheses and models designated as [Adversarial Hypothesis] or [Theoretical Model] are developed within an empirical/simulation environment to test systemic institutional resilience; they do not constitute financial advice, statutory allegations, or institutional endorsements. CERN Zenodo Concept DOI: \href{https://doi.org/10.5281/zenodo.22761112}{10.5281/zenodo.22761112}. Licensed under CC BY 4.0.} \\
    \textit{The Seek Key Research Consortium} \\
    \textit{Working Paper Series in Forensic Epistemology} \\
    \textit{Report No. SK-WP-YYYY-XXX}
}

\date{\today}

\begin{document}

\maketitle

\begin{abstract}
Your abstract here (approx. 150--250 words summarizing problem, methodology, forensic evidence, and theoretical conclusions).
\end{abstract}

\vspace{0.5em}
\noindent \textbf{Keywords:} Keyword 1, Keyword 2, Keyword 3.

\vspace{1em}
\begin{center}
\small
\begin{tabular}{ll}
\toprule
\textbf{Paper Identifier:} & SK-WP-YYYY-XXX \\
\textbf{Evidence Tier:} & \textbf{〔实·账〕} Tier-1 / \textbf{〔推·模〕} Tier-2 / \textbf{〔设·辨〕} Tier-3 \\
\textbf{Subject Classification:} & JEL: G14, G18, D84 \quad | \quad ACM / arXiv: cs.AI, physics.soc-ph \\
\textbf{Permanent DOI:} & \href{https://doi.org/10.5281/zenodo.22761112}{10.5281/zenodo.22761112} (Concept) \\
\bottomrule
\end{tabular}
\end{center}

\newpage

\section{Introduction}

State the fundamental research question and forensic premise.

\section{Epistemological Methodology}

Describe mathematical, physical, or accounting mechanisms.

\section{Forensic Evidence and Derivations}

Present empirical tables, formulas, and balance-sheet proofs.

\section{Conclusion and Discussion}

Summarize findings and policy / capital implications.

\bibliographystyle{plain}
\bibliography{references}

\end{document}
